% Part: incompleteness
% Chapter: introduction
% Section: decidability

\documentclass[../../../include/open-logic-section]{subfiles}

\begin{document}

\olfileid{inc}{int}{dec}

\olsection{అనిర్ణయనీయత మరియు అసంపూర్ణత}

అసంపూర్ణతా సిద్ధాంతాలకు గ్యోడెల్ ఇచ్చిన నిరూపణకు వాక్యనిర్మాణపు అంకగణితీకరణ
అవసరం. కానీ అది లేకుండానే, ఒక సిద్ధాంతం అన్ని
\tetoken{నిర్ణయించదగిన}{!!{decidable}} సంబంధాలకూ
\tetoken{ప్రాతినిధ్యం వహిస్తుందనే}{!!{represents}} పరికల్పన నుండే కొన్ని చక్కని
ఫలితాలు పొందవచ్చు. ఈ నిరూపణ ఆగే సమస్య అనిర్ణయనీయతా నిరూపణను పోలిన కర్ణీయ వాదన.
\sourcecorrection{OLTEINCINT-011}{ఆంగ్ల మూలంలోని ``G\"odel's proof ... require''లో ఏకవచన కర్తకు బహువచన క్రియ వచ్చింది. గ్యోడెల్ నిరూపణకు అంకగణితీకరణ అవసరమనే ఉద్దేశిత అర్థాన్ని ఇచ్చాం.}

\begin{thm}
$\Gamma$ ప్రతి \tetoken{నిర్ణయించదగిన}{!!{decidable}} సంబంధానికీ
\tetoken{ప్రాతినిధ్యం వహించే}{!!{represents}} అవైరుధ్య సిద్ధాంతమైతే, $\Gamma$
\tetoken{నిర్ణయించదగినది}{!!{decidable}} కాదు.
\end{thm}

\begin{proof}
$\Gamma$ \tetoken{నిర్ణయించదగినది}{!!{decidable}} అనుకుందాం. ప్రతి
\tetoken{నిర్ణయించదగిన}{!!{decidable}} సంబంధానికీ $\Gamma$
\tetoken{ప్రాతినిధ్యం వహిస్తే}{!!{represents}}, అది వైరుధ్యమవుతుందని చూపుతాం.

\tetoken{నిర్ణయించదగిన}{!!^{decidable}} లక్షణాలకు (ఏకస్థానిక సంబంధాలకు) ఒక
స్వేచ్ఛా చరం గల \tetoken{సూత్రాలు}{!!{formula}s} ప్రాతినిధ్యం వహిస్తాయి. అటువంటి
అన్ని \tetoken{సూత్రాల}{!!{formula}s} గణనీయ లెక్కింపు $!A_0(x)$,
$!A_1(x)$,~\dots అనుకుందాం. ఇప్పుడు కింది సమితి~$D\subseteq\Nat$ను పరిగణించండి:
\[
D = \Setabs{n}{\Gamma \Proves \lnot !A_n(\num{n})}
\]
$D$ \tetoken{నిర్ణయించదగినది}{!!{decidable}}. ఎందుకంటే $n\in D$ అవుతుందో
మొదట $!A_n(x)$ను, దాని నుండి $\lnot !A_n(\num{n})$ను గణించి పరీక్షించవచ్చు.
$!A_n(x)$లో $x$ స్వేచ్ఛగా వచ్చే ప్రతి చోటా పదం~$\num{n}$ను ప్రతిస్థాపించి,
$!A_n(\num{n})$కు ముందు $\lnot$ చేర్చడం యాంత్రికమైన పని. పరికల్పన ప్రకారం
$\Gamma$ \tetoken{నిర్ణయించదగినది}{!!{decidable}} కనుక,
$\lnot !A_n(\num{n})\in\Gamma$ అవుతుందో పరీక్షించవచ్చు. అయితే $n\in D$;
కాకపోతే $n\notin D$. కాబట్టి $D$ కూడా \tetoken{నిర్ణయించదగినది}{!!{decidable}}.
\sourcecorrection{OLTEINCINT-012}{ఆంగ్ల మూలం $!A_n(x)$ నుండి వికర్ణ వాక్యాన్ని నిర్మిస్తున్న పేరాలో రెండు చోట్ల ఉపలిపి~$n$ను వదిలి $!A(\num{n})$ అని రాసింది. $D$ నిర్వచనానికీ లెక్కింపులోని $n$వ సూత్రానికీ అనుగుణంగా రెండు చోట్లా $!A_n(\num{n})$ను పునరుద్ధరించాం.}

$\Gamma$ అన్ని \tetoken{నిర్ణయించదగిన}{!!{decidable}} లక్షణాలకూ
\tetoken{ప్రాతినిధ్యం వహిస్తుంది}{!!{represents}} కనుక, $D$కు
\tetoken{ప్రాతినిధ్యం వహిస్తుంది}{!!{represents}}. $\Gamma$లో $D$కు
\tetoken{ప్రాతినిధ్యం వహించే}{!!{represents}s}
\tetoken{సూత్రాలన్నీ}{!!{formula}s} $!A_0(x)$, $!A_1(x)$,~\dots మధ్యే ఉన్నాయి.
కాబట్టి $!A_d(x)$ $\Gamma$లో $D$కు
\tetoken{ప్రాతినిధ్యం వహించేలా}{!!{represents}} ఒక సంఖ్య~$d$ను ఎంచుకుందాం.
$d\notin D$ అయితే, $!A_d(x)$ $D$కు
\tetoken{ప్రాతినిధ్యం వహిస్తుంది}{!!{represents}} కనుక,
$\Gamma\Proves\lnot !A_d(\num{d})$. కానీ దాని అర్థం $D$ను నిర్వచించే షరతును
$d$ నెరవేరుస్తుంది; అందువల్ల $d\in D$. ఇది $d\notin D$కు వైరుధ్యం. కనుక
పరోక్ష నిరూపణచే $d\in D$.

$d\in D$ కనుక, $D$ నిర్వచనం ప్రకారం $\Gamma\Proves\lnot !A_d(\num{d})$.
మరోవైపు, $!A_d(x)$ $\Gamma$లో~$D$కు
\tetoken{ప్రాతినిధ్యం వహిస్తుంది}{!!{represents}} కనుక,
$\Gamma\Proves !A_d(\num{d})$. అందువల్ల $\Gamma$ వైరుధ్యమైనది.
\end{proof}

\begin{explain}
అన్ని \tetoken{నిర్ణయించదగిన}{!!{decidable}} సంబంధాలకూ
\tetoken{ప్రాతినిధ్యం వహించే}{!!{represents}}
అవైరుధ్య సిద్ధాంతమేదీ \tetoken{నిర్ణయించదగినది}{!!{decidable}} కాలేదని ముందటి
సిద్ధాంతం చూపుతుంది. అన్ని \tetoken{నిర్ణయించదగిన}{!!{decidable}} సంబంధాలకూ
$\Th{Q}$ నిజంగా \tetoken{ప్రాతినిధ్యం వహిస్తుందని}{!!{represents}s} చూపుతాం;
అందువల్ల $\Th{PA}$, $\Th{TA}$ వంటి $\Th{Q}$ను కలిగి ఉన్న అన్ని సిద్ధాంతాలూ
అలాగే చేస్తాయి, కాబట్టి అవి కూడా \tetoken{నిర్ణయించదగినవి}{!!{decidable}} కావు.
(ఈ సిద్ధాంతాలన్నీ ప్రామాణిక నమూనాలో సత్యమైనవి కనుక అవైరుధ్యమైనవి.)

మొదటి అసంపూర్ణతా సిద్ధాంతపు బలహీన రూపాన్ని పొందడానికి కూడా ఈ ఫలితాన్ని
ఉపయోగించవచ్చు. \tetoken{స్వీకృతీకరించదగిన}{!!{axiomatizable}},
\tetoken{సంపూర్ణమైన}{!!{complete}} ప్రతి సిద్ధాంతమూ
\tetoken{నిర్ణయించదగినది}{!!{decidable}}. కాబట్టి అవైరుధ్యమైన,
\tetoken{స్వీకృతీకరించదగిన}{!!{axiomatizable}}, అన్ని
\tetoken{నిర్ణయించదగిన}{!!{decidable}} లక్షణాలకూ
\tetoken{ప్రాతినిధ్యం వహించే}{!!{represents}s} సిద్ధాంతాలు
\tetoken{సంపూర్ణమైనవి}{!!{complete}} కాలేవు.
\sourcecorrection{OLTEINCINT-013}{ఈ వివరణలోని ఆంగ్ల మూలం ``does represents+s'' మరియు ``represents+s all'' అని రెండుసార్లు తప్పుడు క్రియారూపాన్ని వాడింది. ముడి OpenLogic గుర్తింపు కీలను లక్ష్య టోకెన్లలో యథాతథంగా ఉంచి, రెండు సందర్భాల్లోనూ సరైన ప్రాతినిధ్య అర్థాన్ని ఇచ్చాం.}
\end{explain}

\begin{thm}
$\Gamma$ \tetoken{స్వీకృతీకరించదగినదీ}{!!{axiomatizable}},
\tetoken{సంపూర్ణమైనదీ}{!!{complete}} అయితే, అది
\tetoken{నిర్ణయించదగినది}{!!{decidable}}.
\end{thm}

\begin{proof}
ప్రతి వైరుధ్య సిద్ధాంతమూ \tetoken{నిర్ణయించదగినది}{!!{decidable}}. ఎందుకంటే
వైరుధ్య సిద్ధాంతాలు అన్ని \tetoken{వాక్యాలను}{!!{sentence}s} కలిగి ఉంటాయి;
కాబట్టి ``$!A\in\Gamma$ అవుతుందా?'' అనే ప్రశ్నకు సమాధానం ఎప్పుడూ ``అవును,''
అంటే దాన్ని నిర్ణయించవచ్చు.

ఇప్పుడు $\Gamma$ అవైరుధ్యమనీ, అదనంగా
\tetoken{స్వీకృతీకరించదగినదీ}{!!{axiomatizable}},
\tetoken{సంపూర్ణమైనదీ}{!!{complete}} అనుకుందాం. $\Gamma$
\tetoken{స్వీకృతీకరించదగినది}{!!{axiomatizable}} కనుక అది
\tetoken{గణనీయంగా లెక్కించదగినది}{!!{computably enumerable}}. ఎందుకంటే
$\Gamma$ స్వీకృతాల నుండి వచ్చే సరైన \tetoken{వ్యుత్పత్తులన్నిటినీ}{!!{derivation}s}
ఒక గణనీయ ప్రమేయంతో లెక్కించవచ్చు. సరైన \tetoken{వ్యుత్పత్తి}{!!{derivation}}
నుండి అది \tetoken{నిగమించే}{!!{derive}s} \tetoken{వాక్యాన్ని}{!!{sentence}}
గణించవచ్చు; కాబట్టి ఈ రెండింటినీ కలిపి~$\Gamma$ సిద్ధాంతఫలితాలన్నిటినీ లెక్కించే
గణనీయ ప్రమేయం ఉంటుంది. $\Gamma$ అవైరుధ్యమూ
\tetoken{సంపూర్ణమూ}{!!{complete}} కనుక, \tetoken{వాక్యం}{!!{sentence}}
$\Gamma$ సిద్ధాంతఫలితం అయితేనే $\lnot !A$ సిద్ధాంతఫలితం కాదు. అందువల్ల
$!A\in\Gamma$ అవుతుందో ఇలా నిర్ణయించవచ్చు. $\Gamma$ సిద్ధాంతఫలితాలన్నిటినీ
లెక్కించండి. ఈ జాబితాలో $!A$ కనిపిస్తే $\Gamma\Proves !A$ అని తెలుస్తుంది.
జాబితాలో $\lnot !A$ కనిపిస్తే $\Gamma\Proves/ !A$ అని తెలుస్తుంది. $\Gamma$
\tetoken{సంపూర్ణం}{!!{complete}} కనుక ఈ రెండు సందర్భాల్లో ఒకటి చివరకు సంభవిస్తుంది;
అందువల్ల ప్రక్రియ చివరకు సమాధానం ఇస్తుంది.
\end{proof}

\begin{cor}
\ollabel{cor:incompleteness}
$\Gamma$ అవైరుధ్యమైనదీ, \tetoken{స్వీకృతీకరించదగినదీ}{!!{axiomatizable}}, ప్రతి
\tetoken{నిర్ణయించదగిన}{!!{decidable}} లక్షణానికీ
\tetoken{ప్రాతినిధ్యం వహించేదీ}{!!{represents}} అయితే, అది
\tetoken{సంపూర్ణం}{!!{complete}} కాదు.
\end{cor}

\begin{proof}
$\Gamma$ \tetoken{సంపూర్ణమైతే}{!!{complete}}, అది
\tetoken{స్వీకృతీకరించదగినదీ}{!!{axiomatizable}}, అవైరుధ్యమైనదీ కనుక ముందటి
సిద్ధాంతం ప్రకారం \tetoken{నిర్ణయించదగినది}{!!{decidable}} అవుతుంది. కానీ
$\Gamma$ ప్రతి \tetoken{నిర్ణయించదగిన}{!!{decidable}} లక్షణానికీ
\tetoken{ప్రాతినిధ్యం వహిస్తుంది}{!!{represents}} కనుక, మొదటి సిద్ధాంతం ప్రకారం
అది \tetoken{నిర్ణయించదగినది}{!!{decidable}} కాదు.
\end{proof}

\begin{prob}
$\Th{TA}=\Setabs{!A}{\Sat{N}{!A}}$
\tetoken{స్వీకృతీకరించదగినది}{!!{axiomatizable}} కాదని చూపండి. అన్ని నిర్ణయించదగిన
లక్షణాలకూ $\Th{TA}$ ప్రాతినిధ్యం వహిస్తుందని పరికల్పించవచ్చు.
\end{prob}

ఉదాహరణకు, అన్ని \tetoken{నిర్ణయించదగిన}{!!{decidable}} లక్షణాలకూ $\Th{Q}$
\tetoken{ప్రాతినిధ్యం వహిస్తుందని}{!!{represents}} స్థాపించిన తరువాత, $\Th{Q}$
అసంపూర్ణమని ఉపసిద్ధాంతం చెబుతుంది. అయితే దాని నిరూపణ స్వతంత్ర
\tetoken{వాక్యానికి}{!!{sentence}} ఉదాహరణను ఇవ్వదు; అటువంటి
\tetoken{ఒక వాక్యం}{!!a{sentence}} తప్పకుండా ఉండాలని మాత్రమే చూపుతుంది. ఉదాహరణ
కోసం వాక్యనిర్మాణాన్ని అంకగణితీకరించి గ్యోడెల్ అసలు నిరూపణ ఆలోచనను అనుసరించాలి.
అలాగే, అన్ని \tetoken{నిర్ణయించదగిన}{!!{decidable}} లక్షణాలకూ $\Th{Q}$ నిజంగా
\tetoken{ప్రాతినిధ్యం వహిస్తుందనే}{!!{represents}s} మొదటి వాదనను ఇంకా చూపాలి.

ప్రతి \emph{ఆసక్తికర} సిద్ధాంతమూ అసంపూర్ణమో అనిర్ణయనీయమో కాదని గమనించాలి.
ఆసక్తికరమైన గణిత వాస్తవాలను వర్ణించగలిగేంత బలంగా ఉండి కూడా గ్యోడెల్ ఫలితపు
షరతులను నెరవేర్చని సిద్ధాంతాలు ఎన్నో ఉన్నాయి. ఉదాహరణకు,
$\Th{Pres}=\Setabs{!A\in\Lang{L_{A^+}}}{\Sat{N}{!A}}$, అంటే~$\times$ లేని
అంకగణిత భాషలో ప్రామాణిక నమూనాలో సత్యమైన \tetoken{వాక్యాల}{!!{sentence}s}
సమితి, సంపూర్ణమూ నిర్ణయించదగినదీ. ఈ సిద్ధాంతాన్ని ప్రెస్‌బర్గర్ అంకగణితం అంటారు;
కేవలం $\Obj{0}$, $\prime$,~$+$తో సూత్రీకరించగల సహజ సంఖ్యల సత్యాలన్నిటినీ అది
నిరూపిస్తుంది.\sourcecorrection{OLTEINCINT-014}{ఆంగ్ల మూలంలోని ``satisify'' అక్షరదోషాన్ని ఉద్దేశిత ``satisfy'' అర్థంతో అనువదించాం; గ్యోడెల్ ఫలితపు షరతులను ఈ సిద్ధాంతాలు నెరవేర్చవన్నదే వాక్యపు భావం.}

\end{document}
